\documentclass[noauthor,noinstructornotes]{ximera}

\title{Week 9}
\author{Math 1193 Design Team}

\begin{document}
\begin{abstract}
    Activities for Week 9.
\end{abstract}
\maketitle

\begin{problem}
Consider the set defined by \(x = 1\), or \(x < -2\), or \(x > 4\).

\begin{enumerate}
    \item Draw a number line. Shade the given set.
    \item Write the set in interval notation.
\end{enumerate}
\end{problem}

\begin{problem}
Consider the following functions.

\begin{enumerate}
    \item \(f(x) = 2(x - 2)^2(x + 4)\)
    \begin{enumerate}
        \item Solve the equation \(f(x) = 0\).
        \item What does your answer to the previous part mean about the graph of \(f\)?
    \end{enumerate}
    \item \(g(x) = \frac{2x}{(x - 1)^2(3x + 2)}\)
    \begin{enumerate}
        \item Solve the equation \(g(x) = 0\).
        \item Solve the equation \((x - 1)^2(3x + 2) = 0\).
        \item What do your answers to the previous parts mean about the graph of \(g\)?
    \end{enumerate}
\end{enumerate}

\end{problem}

\begin{problem}
We'll investigate the \(x\)-values where certain functions are positive or negative.

\begin{enumerate}
    \item Draw a number line. Indicate on the number line all \(x\)-values
    for which the function \(x - 2\) is positive and those for which
    it is negative (you can do this by plugging in points and testing
    to see whether the output is positive or negative). Would a graph 
    of the function help? 
    Where does the function change from positive to negative?
    \item Do the same thing for \(x + 3\).
    \item Do the same thing for \((x - 2)(x + 3)\). Use your answers to
    the previous two parts to help.
    \item Organize the information you've found (the functions and where they
    are positive and negative) in a way that works for
    you. Compare your method with a neighbor.
\end{enumerate}

\begin{instructorNotes}
Let students come up with their own organization methods. Make sure to show
a more standard way as a class.
\end{instructorNotes}
\end{problem}

\begin{problem}
Let's play a game! Divide into pairs and choose who will be
Player 1 and Player 2. 

\begin{enumerate}
    \item Player 1 chooses a sequence of 3 operations (addition, 
    multiplication, subtraction or division) in secret. For example:
    \begin{itemize}
        \item Add 2, multiply by 4, subtract 1
    \end{itemize}
    \item Player 1 writes down a function \(f\) which applies those operations,
    in order, to the variable \(x\). For example:
    \begin{itemize}
        \item \(f(x) = 4(x + 2) - 1\)
    \end{itemize}
    \item Player 1 chooses 3 \(x\)-values in secret. Then, they plug them into their function,
    and write down the outputs. For example:
    \begin{itemize}
        \item \(f(0) = 7\), \(f(1) = 11\) and \(f(0.5) = 9\). So Player 1
        writes down 7, 11, 9.
    \end{itemize}
    \item Player 1 shows Player 2 their function \(f\) and the outputs they wrote down.
    For example:
    \begin{itemize}
        \item ``My function is \(f(x) = 4(x + 2) - 1\), and my outputs are 7, 11, and 9.''
    \end{itemize}
    \item Player 2 then tries to determine the original \(x\)-values that
    Player 1 chose. For example:
    \begin{itemize}
        \item ``Your original \(x\)-values were 0, 11, and 0.5.''
    \end{itemize}
    \item Player 2 writes down a sequence of 3 operations that they
    used to determine those \(x\)-values. For example:
    \begin{itemize}
        \item Add 1, divide by 4, subtract 2
    \end{itemize}
    \item Together, the players write down a function \(g\) that always gives the 
    starting value if you're given an output. Base this on your operations
    from the previous part. For example:
    \begin{itemize}
        \item \(g(x) = \frac{x + 1}{4} - 2\)
    \end{itemize}
    \item Take a second together to consider how \(g\) and \(f\) are related. 
    \item Switch roles and play again!
\end{enumerate}
\end{problem}

\end{document}
